\documentclass[onecolumn]{aastex7}

\newcommand{\Header}[1]{{\large{\bf{#1}}}}


\begin{document}

\title{Job Interview Advice}
\author[0000-0002-6134-8946]{Ilya Mandel}
\affiliation{School of Physics and Astronomy, Monash University, Clayton VIC 3800, Australia}
\affiliation{OzGrav: The ARC Centre of Excellence for Gravitational Wave Discovery, Australia}
\email{ilya.mandel@monash.edu}


\begin{abstract}
I attach below some general thoughts on faculty job interviews, though some of these may also be relevant for postdoctoral interviews.
\end{abstract}

\maketitle

\vspace{0.1in}

Whether this works out or not, you’ve done well to be on the list. That means your stock is bullish, and more opportunities will come your way.

Now, what the committee members actually have in mind and what you can do to impress them (if anything) is entirely out of your hands. They may be doing a genuine open search, or they may have an internal candidate they’re going to give the job to no matter what and you might just be filler -- you won’t know and can’t change this. So just do the best you can, and don’t worry about the outcome, which is out of your hands.

Do listen to subtle clues in the interview. If they seem to be begging you to give a particular answer, give it to them (within reason). E.g., if you tell them how excited you are about topic X and they seem bored and are asking whether you might be interested in topic Y, for which they just got a huge institute-wide grant, don’t just keep on talking about X.  But don’t try to sell them some story which you are guessing they want at the expense of being yourself (most likely, they’ll see through it, and be much less impressed than they would have been with your true self).

Also, do make an effort to find out as much as you can about the institute and the people there. It’ll impress the panel if you know who you are talking to and can relate to their work. Conversely, if you haven’t done this, they may think that you aren’t very interested. I was on a search panel where a candidate referred multiple times to a particular paper, without once mentioning a competing paper that just happened to be written by the chair of the department / head of the search panel; let’s just say the chair wasn’t very impressed.  By the way, if the invitation for the interview provides limited information, it is reasonable to ask things like ``Could you tell me who I will be speaking with during the interview'?' or ``At what level should I pitch my presentation?''

The panel are likely asking themselves some of the following questions: Can you describe your work and plans clearly and coherently (i.e., will you have a shot in grant applications)? Are you professional and friendly (i.e., would you be a reasonable colleague)? Are you a decent speaker (i.e., could you teach a course)?

Depending on what they are looking for, their questions for you can range from ``Imagine you had to design a new course on topic X; what subject matter would you put into it, and what innovations would you consider relative to the way you were taught?'' to ``Unfortunately, we all have to do service work along with teaching and research; what would you like to do if you had the choice?'' to ``Where will you apply for grants and what projects would you try to obtain funding for?''

Try to articulate a vision for going forward. What will you be doing in 5 to 10 years? What grants will you apply for? How will you lead the field? What PhD projects would you think of giving students if you are going to have access to PhD students (or undergrad research students, whoever is applicable)? If you will be teaching, think of what class you might particularly like to teach, especially if you have some innovative idea.

Lastly, try to err on the formal side: if you don’t know the people on the panel, assume that they are dolts with a limited sense of humour. 

Perhaps try to think of a suitable final question or set of questions, because they will always ask in the end if you have questions, and this is their last impression of you. Of course, you can ask practical questions like ``How heavy will my teaching load be?'' or ``When will I find out the results?'', but they’re likely to find forward-looking questions about, say, possible future collaborative projects to be more indicative of your excitement and willingness to engage.

Here are some potential questions and general interview advice from other astronomers; I don’t necessarily agree with everything, but these can still be useful to ponder:

\url{https://bryangaensler.net/papers/gaensler_maddison_jobs_2012.pdf}

\url{http://www.astrobetter.com/wiki/Interview+Advice}

\end{document}